\documentclass[a4paper,11pt]{article}
\usepackage[utf8]{inputenc}
\usepackage[T1]{fontenc}
\usepackage[ngerman]{babel}
\usepackage{geometry}
\geometry{left=2.5cm, right=2.5cm, top=2.5cm, bottom=2.5cm}
\usepackage{longtable}
\usepackage{booktabs}
\usepackage{array}
\usepackage{xcolor}
\usepackage{colortbl}
\usepackage{hyperref}
\usepackage{fancyhdr}
\usepackage{listings}
\usepackage{tabularx}
\usepackage{enumitem}

\definecolor{headerblue}{HTML}{1565C0}
\definecolor{tableheader}{HTML}{E3F2FD}
\definecolor{codebg}{HTML}{F5F5F5}
\definecolor{codegreen}{HTML}{2E7D32}
\definecolor{codegray}{HTML}{757575}

\lstset{
  backgroundcolor=\color{codebg},
  basicstyle=\ttfamily\small,
  keywordstyle=\color{headerblue}\bfseries,
  commentstyle=\color{codegreen},
  stringstyle=\color{codegray},
  breaklines=true,
  frame=single,
  rulecolor=\color{codegray!30},
  numbers=none,
  showstringspaces=false,
  tabsize=2,
  language=C++
}

\hypersetup{colorlinks=true,linkcolor=headerblue,urlcolor=headerblue}

\pagestyle{fancy}
\fancyhf{}
\fancyhead[L]{\textbf{INSPECTAIR}}
\fancyhead[R]{Configuration Management v2.0}
\fancyfoot[C]{\thepage}
\renewcommand{\headrulewidth}{0.4pt}

\setlength{\parindent}{0pt}
\setlength{\parskip}{6pt}

\begin{document}

\begin{center}
{\LARGE\textbf{INSPECTAIR}}\\[0.3cm]
{\Large Configuration Management}\\[0.3cm]
{\large Konfigurations- und Versionsverwaltung}\\[0.5cm]
\begin{tabular}{ll}
\textbf{Projekt:} & InspectAir -- Luftqualitätsmessgerät \\
\textbf{Version:} & 2.0 \\
\textbf{Datum:} & Februar 2026 \\
\textbf{Autoren:} & Maximilian Liebl, Cedric Stroh, Jannis Messer \\
\textbf{Modul:} & Mikrocomputertechnik 1, DHBW Ravensburg \\
\textbf{Dozent:} & Thomas Stingl (Airbus) \\
\end{tabular}
\end{center}

\vspace{0.5cm}
\noindent\textit{Dieses Dokument wurde unter Verwendung von KI-Tools (Claude, Anthropic) zur Überprüfung erstellt.}
\vspace{0.3cm}

\tableofcontents
\newpage

% ============================================================
\section{Versionskontrolle}

\begin{tabularx}{\textwidth}{lX}
\toprule
\rowcolor{tableheader}
\textbf{Parameter} & \textbf{Wert} \\
\midrule
System & Git \\
Hosting & GitHub (privates Repository) \\
Branching-Modell & \texttt{main} (stabil) + Feature-Branches \\
Commit-Konvention & Conventional Commits (\texttt{feat:}, \texttt{fix:}, \texttt{docs:}, \texttt{refactor:}) \\
\bottomrule
\end{tabularx}

\subsection{Branch-Strategie}

Entwicklung erfolgt auf Feature-Branches, die nach Review per Pull Request in \texttt{main} gemergt werden. Der \texttt{main}-Branch enthält stets den aktuellen, kompilierbaren Stand. Feature-Branches folgen dem Schema \texttt{feature/<beschreibung>} (z.\,B. \texttt{feature/lvgl9-migration}, \texttt{feature/power-manager}).

% ============================================================
\section{Build-System}

\begin{tabularx}{\textwidth}{lX}
\toprule
\rowcolor{tableheader}
\textbf{Parameter} & \textbf{Wert} \\
\midrule
IDE & VSCode mit PlatformIO Extension \\
Build-System & PlatformIO Core \\
Platform & Espressif32 \\
Board & \texttt{esp32-s3-devkitc-1} \\
Framework & Arduino \\
Flash Mode & QIO 80\,MHz (8\,MB) \\
PSRAM & OPI 80\,MHz (8\,MB) \\
Upload Speed & 921600 Baud \\
Monitor Speed & 115200 Baud \\
\bottomrule
\end{tabularx}

\subsection{Build-Flags}

Zentrale Compiler-Flags aus \texttt{platformio.ini}:

\begin{lstlisting}
build_flags =
  -DARDUINO_USB_CDC_ON_BOOT=1
  -DBOARD_HAS_PSRAM
  -DLV_CONF_INCLUDE_SIMPLE
  -DLV_COLOR_DEPTH=16
  -DLV_COLOR_16_SWAP=1
  -DLV_MEM_SIZE="(96U * 1024U)"
  -DLV_FONT_MONTSERRAT_16=1
\end{lstlisting}

% ============================================================
\section{Bibliotheken}

\begin{tabularx}{\textwidth}{llX}
\toprule
\rowcolor{tableheader}
\textbf{Bibliothek} & \textbf{Version} & \textbf{Verwendung} \\
\midrule
LovyanGFX & 1.1.16 & Display-Treiber für ST7796S \\
LVGL & 9.2.2 & UI-Framework \\
Adafruit AHTX0 & latest & AHT20 Temperatur-/Feuchte-Sensor \\
Adafruit SGP40 & latest & SGP40 VOC-Sensor \\
MH-Z19 & 1.5.4 & MH-Z19C CO\textsubscript{2}-Sensor \\
PMS Library & 1.1.0 & PMS5003 Feinstaub-Sensor \\
ld2410 & 0.1.3 & LD2410C Radar-Sensor \\
EspSoftwareSerial & 8.2.0 & Software-UART für MH-Z19C \\
\bottomrule
\end{tabularx}

% ============================================================
\section{Quellcode-Struktur}

Das Projekt umfasst 30 Quellcode-Dateien:

\begin{tabularx}{\textwidth}{llX}
\toprule
\rowcolor{tableheader}
\textbf{Verzeichnis} & \textbf{Dateien} & \textbf{Beschreibung} \\
\midrule
\texttt{src/} & \texttt{main.cpp} & Hauptprogramm \\
\texttt{include/} & \texttt{pins.h}, \texttt{sensor\_types.h}, \texttt{colors.h} & Zentrale Konfiguration \\
\texttt{include/} & \texttt{display\_config.h}, \texttt{app\_config.h} & Display- und App-Konfiguration \\
\texttt{src/sensors/} & \texttt{aht\_sgp.cpp/h}, \texttt{mhz19c.cpp/h} & I2C- und SoftSerial-Sensoren \\
\texttt{src/sensors/} & \texttt{pms5003.cpp/h}, \texttt{ld2410c.cpp/h} & UART-Sensoren \\
\texttt{src/sensors/} & \texttt{sensor\_filter.cpp/h}, \texttt{sensor\_history.cpp/h} & Datenverarbeitung \\
\texttt{src/display/} & \texttt{lvgl\_driver.cpp/h} & LVGL~9 Display-Treiber \\
\texttt{src/display/} & \texttt{ui\_manager.cpp/h} & Multi-Screen UI (5 Screens) \\
\texttt{src/system/} & \texttt{power\_manager.cpp/h} & Energieverwaltung \\
\texttt{src/network/} & \texttt{WifiClock.cpp/h}, \texttt{web\_remote.cpp/h} & WiFi, NTP, HTTP-Server \\
\bottomrule
\end{tabularx}

% ============================================================
\section{Pin-Konfiguration}

Alle GPIO-Pin-Definitionen sind zentral in \texttt{pins.h} gesammelt und als \texttt{\#define} definiert. Einzelne Module dürfen keine eigenen Pin-Nummern hart codieren, sondern referenzieren stets die Definitionen aus \texttt{pins.h}.

Beispiel:
\begin{lstlisting}
// pins.h - Zentrale Pin-Definitionen
#define TFT_MOSI    11
#define TFT_MISO    13
#define TFT_SCLK    12
#define TFT_CS      14
#define TFT_DC       9
#define TFT_RST     46
#define TFT_BL       3
#define I2C_SDA      8
#define I2C_SCL     18
\end{lstlisting}

% ============================================================
\section{Dokumentenversionen}

\begin{tabularx}{\textwidth}{lll}
\toprule
\rowcolor{tableheader}
\textbf{Dokument} & \textbf{Version} & \textbf{Status} \\
\midrule
Requirements Specification & v2.0 & Aktualisiert \\
Design Description & v2.0 & Aktualisiert \\
Coding Rules & v2.0 & Aktualisiert \\
Configuration Management & v2.0 & Aktualisiert \\
Risk \& Opportunities & v2.0 & Aktualisiert \\
Schedule & v2.0 & Aktualisiert \\
Test Specification & v2.0 & Aktualisiert \\
Pin-Belegung / Bauteilliste & v2.0 & Aktualisiert \\
Requirements Verification & v1.0 & Neu \\
Projektbericht & v1.0 & In Arbeit \\
\bottomrule
\end{tabularx}

% ============================================================
\section{Änderungshistorie}

\begin{tabularx}{\textwidth}{l l l X}
\toprule
\rowcolor{tableheader}
\textbf{Version} & \textbf{Datum} & \textbf{Autor} & \textbf{Änderung} \\
\midrule
1.0 & Januar 2026 & Team & Erstversion \\
2.0 & Februar 2026 & Team & Aktualisierte Bibliotheksversionen (LVGL 9.2.2, LovyanGFX 1.1.16), 30 Quellcode-Dateien, Build-Flags ergänzt, Pin-Konfiguration aus \texttt{pins.h}, Dokumentenversionen aktualisiert. \\
\bottomrule
\end{tabularx}

\end{document}
